% Строчки для суммарной таблицы, в следующей строке прописаны используемые экземпляры узлов
%===s control, LLN0, XCBR1, ESXSWI1, DSXSWI1
\begin{flushleft}
\needspace{4\baselineskip}
\color{uniblue}{\subsection{Коммутационные аппараты (КА)}}
\color{black}
\nopagebreak
\vspace{3mm}
\end{flushleft}

%%%%%%%%%%%%%%%%%%%%%%%%%%%%%%%%%%%%%%%%%%%%%%%%%%%%%%%%%%%%%%%%%%%%% ОБЩИЕ СВЕДЕНИЯ %%%%%%%%%%%%%%%%%%%%%%%%%%%%%%%%%%%%%%%%%%%%%%%%%%%%%%%%%%%%%%

\begin{flushleft}
\color{uniblue}{\subsubsection{Общие сведения}}
\color{black}
\nopagebreak
\end{flushleft}

\par
Функциональный блок <<КА>> предназначен для реализации функций контроля и управления высоковольтными коммутационными аппаратами (заземляющий нож, разъединитель, выключатель) в составе защищаемого присоединения.  

\par
В функциональном блоке <<КА>> предусматривются функции контроля и управления следующими коммутационными аппаратами:
\begin{itemize}
\item выключатель стороны ВН трансформатора;
\item заземляющий нож трансформатора (со стороны ВН);
\item разъединитель стороны ВН трансформатора.
\end{itemize}

\par
Ввод функционального блока <<КА>> в работу на этапе параметрирования устройства осуществляется программным переключателем SGF1 <<Ввод\_функции>>. Введенное состояние функционального блока соответствует активному состоянию сигнала <<КА: Ввод>>. Функциональный блок <<КА>> может быть выведен из работы оперативно путем активации сигнала <<ОВ КА>>, о чем сигнализирует активное состояние сигнала <<КА: Оперативный вывод>>.

Логика управления режимами работы функционального блока <<КА>> выполнена в соответствии с рисунком \ref{ka:img_gen}. В таблице \ref{ka:tbl_gen} приведены параметры, необходимые для настройки функционального блока <<КА>>. 

\vspace{2mm}
\begin{figure}[H]
\centering

\includegraphics[width=0.8\textwidth,height=0.8\textheight,keepaspectratio]{\fbpath/39. КА/_latex/img/1.pdf}
\captionof{figure}{Логика управления режимами функционального блока <<КА>>}\label{ka:img_gen}
\end{figure}

\begin{longtable}{|>{\centering\arraybackslash}m{5cm}|>{\centering\arraybackslash}m{4cm}|>{\centering\arraybackslash}m{4cm}|>{\centering\arraybackslash}m{2cm}|>{\centering\arraybackslash}m{1cm}|}
\caption{Параметры для настройки функционального блока <<КА>>\hfill\vspace{-0.5\baselineskip}}\label{ka:tbl_gen}\\ 
\hline
\rowcolor{gray!30}
Параметр (Параметр на ИЧМ) & Условное обозначение на схеме & Значение/ Диапазон & Единица измерения & Шаг \\ 
\hline
\endfirsthead
\caption{\emph{Продолжение\hfill\vspace{-0.5\baselineskip}}} \\ % Отступ обозначения таблицы от таблицы и выравнивание надписи слева
\hline
\rowcolor{gray!30}
Параметр (Параметр на ИЧМ) & Условное обозначение на схеме & Значение/ Диапазон & Единица измерения & Шаг \\ 
%\hline
\endhead
%\hline
\endfoot
%\hline
\endlastfoot
%===t1*LLN0
\centering Ввод функции в работу (Ввод\_функции) & \centering SGF1 & \centering 0 = Не предусмотрено\\1 = Предусмотрено & \centering -- & \centering\arraybackslash -- \\
\hline
%===t1
\end{longtable}

%%%%%%%%%%%%%%%%%%%%%%%%%%%%%%%%%%%%%%%%%%%%%%%%%%%%%%%%%%% ВЫКЛЮЧАТЕЛЬ %%%%%%%%%%%%%%%%%%%%%%%%%%%%%%%%%%%%%%%%%%%%%%%%%%%%%%%%%%

\begin{flushleft}
\color{uniblue}{\subsubsection{Выключатель (В)}}
\color{black}
\nopagebreak
\end{flushleft}

\par
Ввод функции <<Выключатель>> в работу на этапе параметрирования устройства осуществляется программным переключателем SGF1 <<Ввод\_функции>>. Введенное состояние функции соответствует активному состоянию сигнала <<В: Ввод>>.

\par
При приеме входных сигналов от выходной логики отключения выключателя стороны ВН трансформатора (<<ЛО: Отключить (Tr)>> - аварийное отключение) или от логики оперативного управления (<<УВ: Отключить>>) формируется команда <<В: Отключить (реле)>> на внутреннее выходное реле устройства с целью коммутации выходной цепи отключения ЭМО выключателя. При этом возможна выдача как импульсного сигнала (режим <<Импульсный>> программного переключателя SGF2 <<Режим\_откл>>, время импульса определяется уставкой Т2 <<Тимп\_откл>>), так и длительного сигнала (режим <<Длительный>> программного переключателя SGF2 <<Режим\_откл>>). 
В режиме выдачи импульсного сигнала выдаваемый импульс (а также и выдаваемый сигнал в длительном режиме) ограничивается активизацией сигнала отключенного состояния выключателя <<В отключен>> с учетом или без учета сигналов <<Работа ЭМО1>> и <<Работа ЭМО2>> (в зависимости от положения программного переключателя SGF3 <<Контр\_раб\_ЭМО>>).

\par
При приеме входного сигнала от логики оперативного управления (<<УВ: Включить>>) формируется команда <<В: Включить (реле)>> на внутреннее выходное реле устройства с целью коммутации выходной цепи включения ЭМВ выключателя. При этом возможна выдача как импульсного сигнала (режим <<Импульсный>> программного переключателя SGF4 <<Режим\_вкл>>, время импульса определяется уставкой Т3 <<Тимп\_вкл>>), так и длительного сигнала (режим <<Длительный>> программного переключателя SGF4 <<Режим\_вкл>>).
В режиме выдачи импульсного сигнала выдаваемый импульс (а также и выдаваемый сигнал в длительном режиме) ограничивается активизацией сигнала включенного состояния выключателя <<В включен>> с учетом или без учета сигнала <<Работа ЭМВ>> (в зависимости от положения программного переключателя SGF5 <<Контр\_раб\_ЭМВ>>). При этом возможна задержка сигнала <<В включен>> (определяется уставкой Т4 <<Тпрод\_вкл>>) для сброса импульсного или длительного сигнала на включение (для исключения ситуаций так называемого <<опрокидывания>> выключателя при раннем съёме команды <<В: Включить (реле)>> -- это характерно для некоторых видов масляных выключателей).

\par
В логике функции <<Выключатель>> реализована блокировка на прохождение команды включения выключателя <<В: Включить (реле)>> при приеме сигнала аварийного отключения от защит.

\par
Логика функции <<Выключатель>> формирует сигналы состояния выключателя <<В: Отключен>>, <<В: Включен>>, <<В: Промеж. положение>> и <<В: неиспр. положение>>, при этом сигнал <<В: неиспр. положение>> выдается с выдержкой времени, задаваемой уставкой Т1 <<Тнеиспр\_В>>).

\par
В логике функции <<Выключатель>> также реализован счетчик количества выполненных операций <<Включить -- Отключить>>, счетчик сбрасывается при активизации входного сигнала <<Сброс счет. опер. В>>. 

\par
Алгоритм работы функции <<Выключатель>> выполнен в соответствии с рисунком \ref{ka:img_v}. В таблице \ref{ka:tbl_v} приведены параметры, необходимые для настройки функционального блока <<Выключатель>>.

\vspace{2mm}
\begin{figure}[H]
\centering
\includegraphics[width=0.8\textwidth,height=0.8\textheight,keepaspectratio]{\fbpath/39. КА/_latex/img/v.pdf}
\captionof{figure}{Функциональная схема <<Выключатель>>}\label{ka:img_v}
\end{figure}

\begin{longtable}{|>{\centering\arraybackslash}m{5cm}|>{\centering\arraybackslash}m{4cm}|>{\centering\arraybackslash}m{4cm}|>{\centering\arraybackslash}m{2cm}|>{\centering\arraybackslash}m{1cm}|}
    \caption{Параметры для настройки функции <<Выключатель>>\hfill\vspace{-0.5\baselineskip}}\label{ka:tbl_v}\\ 
    \hline
    \rowcolor{gray!30}
    Параметр (Параметр на ИЧМ) & Условное обозначение на схеме & Значение/ Диапазон & Единица измерения & Шаг \\ 
    \hline
    \endfirsthead
    \caption{\emph{Продолжение\hfill\vspace{-0.5\baselineskip}}} \\ % Отступ обозначения таблицы от таблицы и выравнивание надписи слева
    \hline
    \rowcolor{gray!30}
    Параметр (Параметр на ИЧМ) & Условное обозначение на схеме & Значение/ Диапазон & Единица измерения & Шаг \\ 
    %\hline
    \endhead
    %\hline
    \endfoot
    %\hline
    \endlastfoot
    %===t1*XCBR1
\centering Ввод функции в работу (Ввод\_функции) & \centering SGF1 & \centering 0 = Не предусмотрено\\1 = Предусмотрено & \centering -- & \centering\arraybackslash -- \\
\hline
\centering Режим отключения В (Режим\_откл) & \centering SGF2 & \centering 0 = Длительный\\1 = Импульсный & \centering -- & \centering\arraybackslash -- \\
\hline
\centering Контроль работы ЭМО (Контр\_раб\_ЭМО) & \centering SGF3 & \centering 0 = Не предусмотрено\\1 = Предусмотрено & \centering -- & \centering\arraybackslash -- \\
\hline
\centering  Режим включения В (Режим\_вкл) & \centering SGF4 & \centering 0 = Длительный\\1 = Импульсный & \centering -- & \centering\arraybackslash -- \\
\hline
\centering Контроль работы ЭМВ (Контр\_раб\_ЭМВ) & \centering SGF5 & \centering 0 = Не предусмотрено\\1 = Предусмотрено & \centering -- & \centering\arraybackslash -- \\
\hline
\centering Длительность формирования сигнала о неисправном положении В (Tнеиспр\_В) & \centering T1 & \centering 0,01 ... 10,00 & \centering с & \centering\arraybackslash 0,01 \\
\hline
\centering Длительность формирования команды <<Отключить В (реле)>> (Тимп\_откл) & \centering T2 & \centering 0,05 ... 10,00 & \centering с & \centering\arraybackslash 0,01 \\
\hline
\centering Длительность формирования команды <<Включить В (реле)>> (Тимп\_вкл) & \centering T3 & \centering 0,05 ... 10,00 & \centering с & \centering\arraybackslash 0,01 \\
\hline
\centering Время для исключения <<опрокидывания>>  В (Тпрод\_вкл) & \centering T4 & \centering 0,00 ... 1,00 & \centering с & \centering\arraybackslash 0,01 \\
\hline
    %===t1
    \end{longtable}

%%%%%%%%%%%%%%%%%%%%%%%%%%%%%%%%%%%%%%%%%%%%%%%%%%%%%%%%%%%%%%% ЗАЗЕМЛЯЮЩИЙ НОЖ %%%%%%%%%%%%%%%%%%%%%%%%%%%%%%%%%%%%%%%%%%%%%%%%%%%%%%%%%%

\begin{flushleft}
\color{uniblue}{\subsubsection{Заземляющий нож (ЗН)}}
\color{black}
\nopagebreak
\end{flushleft}

\par
Ввод функции <<Заземляющий нож>> в работу на этапе параметрирования устройства осуществляется программным переключателем SGF1 <<Ввод\_функции>>. Введенное состояние функции соответствует активному состоянию сигнала <<ЗН: Ввод>>.

\par
При приеме входного сигнала от логики оперативного управления (<<УЗН: Отключить>>) формируется команда <<ЗН: Отключить (реле)>> на внутреннее выходное реле устройства с целью коммутации выходной цепи отключения заземляющего ножа. При этом возможна выдача как импульсного сигнала (режим <<Импульсный>> программного переключателя SGF2 <<Режим\_откл>>, время импульса определяется уставкой Т2 <<Тимп\_откл>>), так и длительного сигнала (режим <<Длительный>> программного переключателя SGF2 <<Режим\_откл>>). 
В режиме выдачи импульсного сигнала выдаваемый импульс (а также и выдаваемый сигнал в длительном режиме) ограничивается активизацией сигнала отключенного состояния заземляющего ножа <<ЗН отключен>>.

\par
При приеме входного сигнала от логики оперативного управления (<<УЗН: Включить>>) формируется команда <<ЗН: Включить (реле)>> на внутреннее выходное реле устройства с целью коммутации выходной цепи включения заземляющего ножа. При этом возможна выдача как импульсного сигнала (режим <<Импульсный>> программного переключателя SGF3 <<Режим\_вкл>>, время импульса определяется уставкой Т3 <<Тимп\_вкл>>), так и длительного сигнала (режим <<Длительный>> программного переключателя SGF3 <<Режим\_вкл>>).
В режиме выдачи импульсного сигнала выдаваемый импульс (а также и выдаваемый сигнал в длительном режиме) ограничивается активизацией сигнала включенного состояния заземляющего ножа <<ЗН включен>>.

\par
Логика функции <<Заземляющий нож>> формирует сигналы состояния выключателя <<ЗН: Отключен>>, <<ЗН: Включен>>, <<ЗН: Промеж. положение>> и <<ЗН: неиспр. положение>>.

\par
В логике функции <<Заземляющий нож>> также реализован счетчик количества выполненных операций <<Включить -- Отключить>>, счетчик сбрасывается при активизации входного сигнала <<Сброс счет. опер. ЗН>>. 

\par
Алгоритм работы функции <<Заземляющий нож>> выполнен в соответствии с рисунком \ref{ka:img_zn}. В таблице \ref{ka:tbl_zn} приведены параметры, необходимые для настройки функционального блока <<Заземляющий нож>>.

\vspace{1mm}
\begin{figure}[H]
\centering
\includegraphics[width=0.75\textwidth,height=0.75\textheight,keepaspectratio]{\fbpath/39. КА/_latex/img/zn.pdf}
\captionof{figure}{Функциональная схема <<ЗН>>}\label{ka:img_zn}
\end{figure}

\begin{longtable}{|>{\centering\arraybackslash}m{5cm}|>{\centering\arraybackslash}m{4cm}|>{\centering\arraybackslash}m{4cm}|>{\centering\arraybackslash}m{2cm}|>{\centering\arraybackslash}m{1cm}|}
    \caption{Параметры для настройки функции <<Заземляющий нож>>\hfill\vspace{-0.5\baselineskip}}\label{ka:tbl_zn}\\ 
    \hline
    \rowcolor{gray!30}
    Параметр (Параметр на ИЧМ) & Условное обозначение на схеме & Значение/ Диапазон & Единица измерения & Шаг \\ 
    \hline
    \endfirsthead
    \caption{\emph{Продолжение\hfill\vspace{-0.5\baselineskip}}} \\ % Отступ обозначения таблицы от таблицы и выравнивание надписи слева
    \hline
    \rowcolor{gray!30}
    Параметр (Параметр на ИЧМ) & Условное обозначение на схеме & Значение/ Диапазон & Единица измерения & Шаг \\ 
    %\hline
    \endhead
    %\hline
    \endfoot
    %\hline
    \endlastfoot
    %===t1*ESXSWI1
\centering Ввод функции в работу (Ввод\_функции) & \centering SGF1 & \centering 0 = Не предусмотрено\\1 = Предусмотрено & \centering -- & \centering\arraybackslash -- \\
\hline
\centering Режим отключения ЗН (Режим\_откл) & \centering SGF2 & \centering 0 = Длительный\\1 = Импульсный & \centering -- & \centering\arraybackslash -- \\
\hline
\centering Режим включения ЗН (Режим\_вкл) & \centering SGF3 & \centering 0 = Длительный\\1 = Импульсный & \centering -- & \centering\arraybackslash -- \\
\hline
\centering Длительность формирования сигнала о неисправном положении ЗН (Tнеиспр\_ЗН) & \centering T1 & \centering 0,01 ... 10,00 & \centering с & \centering\arraybackslash 0,01 \\
\hline
\centering Длительность формирования команды <<Отключить ЗН (реле)>> (Тимп\_откл) & \centering T2 & \centering 0,10 ... 100,00 & \centering с & \centering\arraybackslash 0,01 \\
\hline
\centering Длительность формирования команды <<Включить ЗН (реле)>> (Тимп\_вкл) & \centering T3 & \centering 0,10 ... 100,00 & \centering с & \centering\arraybackslash 0,01 \\
\hline
    %===t1
    \end{longtable}

%%%%%%%%%%%%%%%%%%%%%%%%%%%%%%%%%%%%%%%%%%%%%%%%%%%%%%%%%%%%%%%%%% РАЗЪЕДИНИТЕЛЬ  %%%%%%%%%%%%%%%%%%%%%%%%%%%%%%%%%%%%%%%%%%%%%%%%%%

\begin{flushleft}
\color{uniblue}{\subsubsection{Разъединитель (Р)}}
\color{black}
\nopagebreak
\end{flushleft}

\par
Ввод функции <<Разъединитель>> в работу на этапе параметрирования устройства осуществляется программным переключателем SGF1 <<Ввод\_функции>>. Введенное состояние функции соответствует активному состоянию сигнала <<Р: Ввод>>.

\par
При приеме входного сигнала от логики оперативного управления (<<УР: Отключить>>) формируется команда <<Р: Отключить (реле)>> на внутреннее выходное реле устройства с целью коммутации выходной цепи отключения разъединителя. При этом возможна выдача как импульсного сигнала (режим <<Импульсный>> программного переключателя SGF2 <<Режим\_откл>>, время импульса определяется уставкой Т2 <<Тимп\_откл>>), так и длительного сигнала (режим <<Длительный>> программного переключателя SGF2 <<Режим\_откл>>). 
В режиме выдачи импульсного сигнала выдаваемый импульс (а также и выдаваемый сигнал в длительном режиме) ограничивается активизацией сигнала отключенного состояния разъединителя <<ЗН отключен>>.

\par
При приеме входного сигнала от логики оперативного управления (<<УР: Включить>>) формируется команда <<Р: Включить (реле)>> на внутреннее выходное реле устройства с целью коммутации выходной цепи включения разъединителя. При этом возможна выдача как импульсного сигнала (режим <<Импульсный>> программного переключателя SGF3 <<Режим\_вкл>>, время импульса определяется уставкой Т3 <<Тимп\_вкл>>), так и длительного сигнала (режим <<Длительный>> программного переключателя SGF3 <<Режим\_вкл>>).
В режиме выдачи импульсного сигнала выдаваемый импульс (а также и выдаваемый сигнал в длительном режиме) ограничивается активизацией сигнала включенного состояния разъединителя <<ЗН включен>>.

\par
Логика функции <<Разъединитель>> формирует сигналы состояния выключателя <<Р: Отключен>>, <<Р: Включен>>, <<Р: Промеж. положение>> и <<Р: неиспр. положение>>.

\par
В логике функции <<Разъединитель>> также реализован счетчик количества выполненных операций <<Включить -- Отключить>>, счетчик сбрасывается при активизации входного сигнала <<Сброс счет. опер. Р>>. 

\par
Алгоритм работы функции <<Разъединитель>> выполнен в соответствии с рисунком \ref{ka:img_r}. В таблице \ref{ka:tbl_r} приведены параметры, необходимые для настройки функционального блока <<Разъединитель>>.

\vspace{1mm}
\begin{figure}[H]
\centering
\includegraphics[width=0.75\textwidth,height=0.75\textheight,keepaspectratio]{\fbpath/39. КА/_latex/img/r.pdf}
\captionof{figure}{Функциональная схема <<Р>>}\label{ka:img_r}
\end{figure}

\begin{longtable}{|>{\centering\arraybackslash}m{5cm}|>{\centering\arraybackslash}m{4cm}|>{\centering\arraybackslash}m{4cm}|>{\centering\arraybackslash}m{2cm}|>{\centering\arraybackslash}m{1cm}|}
    \caption{Параметры для настройки функции <<Разъединитель>>\hfill\vspace{-0.5\baselineskip}}\label{ka:tbl_zn}\\ 
    \hline
    \rowcolor{gray!30}
    Параметр (Параметр на ИЧМ) & Условное обозначение на схеме & Значение/ Диапазон & Единица измерения & Шаг \\ 
    \hline
    \endfirsthead
    \caption{\emph{Продолжение\hfill\vspace{-0.5\baselineskip}}} \\ % Отступ обозначения таблицы от таблицы и выравнивание надписи слева
    \hline
    \rowcolor{gray!30}
    Параметр (Параметр на ИЧМ) & Условное обозначение на схеме & Значение/ Диапазон & Единица измерения & Шаг \\ 
    %\hline
    \endhead
    %\hline
    \endfoot
    %\hline
    \endlastfoot
    %===t1*DSXSWI1
\centering Ввод функции в работу (Ввод\_функции) & \centering SGF1 & \centering 0 = Не предусмотрено\\1 = Предусмотрено & \centering -- & \centering\arraybackslash -- \\
\hline
\centering Режим отключения Р (Режим\_откл) & \centering SGF2 & \centering 0 = Длительный\\1 = Импульсный & \centering -- & \centering\arraybackslash -- \\
\hline
\centering Режим включения Р (Режим\_вкл) & \centering SGF3 & \centering 0 = Длительный\\1 = Импульсный & \centering -- & \centering\arraybackslash -- \\
\hline
\centering Длительность формирования сигнала о неисправном положении Р (Tнеиспр\_Р) & \centering T1 & \centering 0,01 ... 10,00 & \centering с & \centering\arraybackslash 0,01 \\
\hline
\centering Длительность формирования команды <<Отключить Р (реле)>> (Тимп\_откл) & \centering T2 & \centering 0,10 ... 100,00 & \centering с & \centering\arraybackslash 0,01 \\
\hline
\centering Длительность формирования команды <<Включить Р (реле)>> (Тимп\_вкл) & \centering T3 & \centering 0,10 ... 100,00 & \centering с & \centering\arraybackslash 0,01 \\
\hline
    %===t1
    \end{longtable}